% Part: propositional-logic

\documentclass[../../include/open-logic-part]{subfiles}

\begin{document}

\olpart{pl}{قضیاتی منطق}

\begin{editorial}
  ایس حصے وچ کلاسیکی قضیاتی منطق دا مواد شامل اے۔ پہلا باب نسبتاً
  ابتدائی اے تے صرف تعریفاں تے نتیجے درج کردا اے؛ کئی ثبوت پورے نہیں
  دِتے گئے بلکہ مشقاں لئی چھڈّے گئے نیں۔ ثبوتی نظاماں تے تمامیت دے
  قضیے دا مواد پہلے درجے دی منطق والے حصے توں شامل کیتا گیا اے، تے
  ``FOL'' ٹیگ نوں false مقرر کیتا گیا اے۔ ایس طرح محمولاں، اصطلاحاں
  تے مسوّراں نال متعلق ہر گل شامل نہیں رہندی، تے
  !!{structure}s~$\Struct{M}$ دی گل نوں
  !!{valuation}s~$\pAssign{v}$ دی گل نال بدل دِتا جاندا اے۔

  منصوبہ اے کہ ایس حصے نوں ہور تفصیل نال ودھایا جاوے، تے ہور موضوع
  تے نتیجے، مثلاً صداقتی-فنکشنی تمامیت، شامل کیتے جان۔
\end{editorial}

\olimport[syntax-and-semantics]{syntax-and-semantics}

\tagfalse{FOL}

\olimport[../first-order-logic/proof-systems]{proof-systems}

\iftag{prfSC}{%
  \olimport[../first-order-logic/sequent-calculus]{sequent-calculus}
}{}

\iftag{prfND}{%
  \olimport[../first-order-logic/natural-deduction]{natural-deduction}
}{}

\iftag{prfTab}{%
  \olimport[../first-order-logic/tableaux]{tableaux}
}{}

\iftag{prfAX}{%
  \olimport[../first-order-logic/axiomatic-deduction]{axiomatic-deduction}
}{}

\olimport[../first-order-logic/completeness]{completeness}

\tagtrue{FOL}

\OLEndPartHook

\end{document}
